Queremos saber cómo calificaría el modelo a vacas que nunca vio. Para eso separamos a las vacas, no a las fotos. Todas las fotos de una vaca están siempre del mismo lado.

\subsection{Partición}

\begin{itemize}
  \item 76 vacas con BCS repartidas en 5 pliegues fijos por vaca (16/15/15/15/15), de la columna \texttt{fold} del índice. Las vacas se ordenan por BCS e id y se reparten en round-robin. Semilla principal 2026.
  \item En cada rotación, un pliegue es la validación externa. Solo mide, una sola vez, al final. Los demás pliegues forman el entrenamiento, según la tabla.
  \item Juntando las 5 validaciones externas se obtiene una calificación para cada una de las 76 vacas; de ahí sale el n/76.
\end{itemize}

\begin{table}[H]
  \centering
  \small
  \caption{Papel de cada pliegue en las cinco rotaciones.}
  \label{tab:rotaciones}
  \begin{tabularx}{\textwidth}{>{\raggedright\arraybackslash}p{0.12\textwidth}>{\raggedright\arraybackslash}X>{\raggedright\arraybackslash}X>{\raggedright\arraybackslash}X}
    \toprule
    Rotación $k$ & No se usa (prueba original) & Validación externa (solo mide) & Entrenamiento \\
    \midrule
    0 & 0 & 1 & 2, 3, 4 \\
    1 & 1 & 2 & 0, 3, 4 \\
    2 & 2 & 3 & 0, 1, 4 \\
    3 & 3 & 4 & 0, 1, 2 \\
    4 & 4 & 0 & 1, 2, 3 \\
    \bottomrule
  \end{tabularx}
\end{table}

Al inicio se trabajó con una copia de 66 vacas. Cuando el bucket de datos se actualizó el 5 de noviembre, todas las cifras finales de este documento se recalcularon con las 76 vacas. Las cifras sobre 66 vacas ya no se usan.

Todas las decisiones (receta de la cabeza, regularización, estiramiento de la salida, tasa de aprendizaje, época de paro, pesos y candidato) se toman con validación cruzada interna, solo dentro de los pliegues de entrenamiento. Gana la opción que deja más vacas del entrenamiento dentro de $\pm0.25$ (desempate: menor MAE y opción más simple). Luego se reajusta con todo el entrenamiento y predice la validación una sola vez.

\subsection{Métricas}

La métrica principal es el porcentaje de vacas dentro de $\pm0.25$ de la calificación del experto, reportado como n/76 (una vaca = 1.5 puntos porcentuales). Nuestra meta principal es acertar el BCS de al menos 70\,\% de las vacas, es decir, 54/76.

También ocupamos las siguientes métricas como apoyo:

\begin{itemize}
  \item MAE por vaca y porcentaje de vacas dentro de $\pm0.50$ (meta: al menos 90\,\%).
  \item Desglose por extremos (vacas de BCS bajo, centro y BCS alto) para ver dónde falla el modelo.
  \item Dos lecturas, siempre juntas. Salida continua y salida redondeada a 0.25.
\end{itemize}

\subsection{¿Por qué se redondea la salida a 0.25?}

El BCS de vacas Holstein se califica en una escala de 1 a 5 con incrementos de 0.25. \textcite{edmonson1989} describe cada zona del cuerpo en esa escala de cuartos de punto, y \textcite{ferguson1994} confirman que los evaluadores asignan el BCS en pasos de 0.25. Un experto nunca emite un valor intermedio como 3.13, emite 3.00 o 3.25. Por eso el sistema entrega la calificación en la misma escala.

Los estudios de calificación automática del BCS reportan sus resultados de la misma forma, como el porcentaje de vacas con coincidencia exacta, dentro de 0.25 y dentro de 0.50 de la calificación del experto: \textcite{yukun2019} y \textcite{nagy2023}, que además trabajan directamente con clases de cuarto de punto.

La regla no tiene parámetros que ajustar y se aplica igual a todos los modelos. Pero redondear amplía la ventana efectiva de acierto. Para una vaca calificada en la rejilla de cuartos de punto, una predicción continua a menos de unos 0.375 del valor real termina contando como acierto. Por eso reportamos siempre las dos cifras, la continua y la redondeada.
